% !TeX program = xelatex
\documentclass[12pt,a4paper]{article}
\usepackage{amsmath,amssymb}
\usepackage{geometry}
\usepackage{booktabs}
\usepackage{xcolor}
\usepackage{tcolorbox}
\usepackage[UTF8]{ctex}
\geometry{margin=2.5cm}

\newtcolorbox{keypoint}[1][]{
  colback=red!5!white, colframe=red!60!black, fonttitle=\bfseries,
  title=#1, sharp corners=south, boxrule=0.8pt
}

\title{\textbf{相线（Phase Line）分析详解}\\[5pt]
\large ——逐因子判断 $f(x)$ 的符号}
\author{}
\date{}

\begin{document}
\maketitle

\section{问题}

模型为：
\begin{equation}\label{eq:model}
f(x) = r \cdot x \cdot \left(1 - \frac{x}{K_{\text{eff}}}\right), \qquad r > 0,\; K_{\text{eff}} > 0
\end{equation}

$x$ 的物理定义域为 $[x_0, +\infty)$（回收量从 $x_0 = 6314$ 出发，$f(x) > 0$ 保证 $x$ 单调递增，不会回落）。
代数方程 $f(x) = 0$ 有两个根：$x = 0$ 和 $x = K_{\text{eff}}$，但 $x = 0$ 远在定义域之外，
物理上只存在一个平衡点：$x^* = K_{\text{eff}}$。

\section{澄清：定义域内只有一个平衡点，它把定义域切成两个区间}

$x^* = K_{\text{eff}} = 8142.22$，初值 $x_0 = 6314$，$x_0 < K_{\text{eff}}$，平衡点将定义域切为两个区间：

\[
\underbrace{\qquad [x_0,\; K_{\text{eff}}) \qquad}_{\text{区间 I}}
\;\; \underbrace{\rule{30pt}{0pt}}_{x^* = K_{\text{eff}}}
\;\; \underbrace{\qquad (K_{\text{eff}},\; +\infty) \qquad}_{\text{区间 II}}
\]

（数学上的全实轴还有 $(-\infty, 0)$ 和 $(0, x_0)$，但它们不在 $x$ 的物理定义域内。）

\section{逐区间、逐因子判断 $f(x)$ 的符号}

$f(x) = r \cdot x \cdot (1 - x/K_{\text{eff}})$ 是\textbf{三个因子的乘积}。
$r = \gamma/(1+\beta) > 0$ 恒成立（$\gamma > 0$，$1+\beta \geq 1$）。

\subsection*{区间 I：$x_0 \leq x < K_{\text{eff}}$}

\begin{center}
\begin{tabular}{lcc|l}
\toprule
因子 & 表达式 & 符号 & 理由 \\
\midrule
因子 1 & $r$                              & $+$ & $\gamma>0,\;1+\beta>0$ 恒成立 \\
因子 2 & $x$                              & $+$ & $x \geq x_0 = 6314 > 0$ \\
因子 3 & $1 - x/K_{\text{eff}}$           & $+$ & $x < K_{\text{eff}} \;\Rightarrow\; x/K_{\text{eff}} < 1$ \\
\midrule
\multicolumn{4}{c}{\textbf{乘积 $f(x) \;=\; \boxed{+}$（正）}} \\
\bottomrule
\end{tabular}
\end{center}

\[
f(x) > 0 \;\Longrightarrow\; \frac{dx}{dt} > 0 \;\Longrightarrow\; x \text{ 随时间 }\textbf{增大}\;\Longrightarrow\; \text{箭头向 }\textbf{右}\;\longrightarrow
\]

\subsection*{区间 II：$x > K_{\text{eff}}$}

\begin{center}
\begin{tabular}{lcc|l}
\toprule
因子 & 表达式 & 符号 & 理由 \\
\midrule
因子 1 & $r$                              & $+$ & $\gamma>0,\;1+\beta>0$ 恒成立 \\
因子 2 & $x$                              & $+$ & $x > K_{\text{eff}} > 0$ \\
因子 3 & $1 - x/K_{\text{eff}}$           & $-$ & $x > K_{\text{eff}} \;\Rightarrow\; x/K_{\text{eff}} > 1$ \\
\midrule
\multicolumn{4}{c}{\textbf{乘积 $f(x) \;=\; \boxed{-}$（负）}} \\
\bottomrule
\end{tabular}
\end{center}

\[
f(x) < 0 \;\Longrightarrow\; \frac{dx}{dt} < 0 \;\Longrightarrow\; x \text{ 随时间 }\textbf{减小}\;\Longrightarrow\; \text{箭头向 }\textbf{左}\;\longleftarrow
\]

\section{画出相线}

物理定义域内的相线：

\[
\begin{array}{cccccc}
& f(x) > 0 & & f(x) < 0 & \\
& \longrightarrow & & \longleftarrow & \\
\hline
x_0 & & x^* = K_{\text{eff}} & & +\infty \\[6pt]
& & \text{渐近稳定} & &
\end{array}
\]

\section{判读}

\begin{keypoint}[稳定性判据（相线版）]
\begin{itemize}
  \item 两边箭头\textbf{都指向}平衡点 $\rightarrow$ \textbf{稳定}（汇，sink）
  \item 两边箭头\textbf{都背离}平衡点 $\rightarrow$ \textbf{不稳定}（源，source）
  \item 一边指向、一边背离 $\rightarrow$ 半稳定（鞍点退化）
\end{itemize}
\end{keypoint}

在物理定义域内只有一个平衡点 $x^* = K_{\text{eff}}$，两边箭头都指向它：
区间 I 中 $f(x) > 0$（$\rightarrow$），区间 II 中 $f(x) < 0$（$\leftarrow$）
$\Rightarrow$ \textbf{渐近稳定}。

\section{物理直觉}

\begin{itemize}
  \item 回收量在 $[x_0, K_{\text{eff}})$ 之间 $\rightarrow$ 还有剩余增长空间，系统\textbf{自动增长}
  \item 回收量超过 $K_{\text{eff}}$（例如受外部冲击，如年末大扫除造成短期暴增）$\rightarrow$ 超出承载能力，系统\textbf{自动回落}
  \item \textbf{从定义域内任何一点出发，最终都稳定在 $K_{\text{eff}}$}
\end{itemize}

（数学上 $x = 0$ 也是根且不稳定——这恰好解释了为什么回收体系不会自发``死掉''，但 $x = 0$ 在定义域 $[x_0, +\infty)$ 之外，不是系统的有效平衡点。）

\section{相线与线性化的等价关系}

线性化方法算的是 $f'(x^*)$ 的符号，相线方法看的是 $f(x)$ 在平衡点两侧的符号。
本模型中 $f'(x) = r(1 - 2x/K_{\text{eff}})$，
$f'(K_{\text{eff}}) = -r < 0$（稳定），与相线结论完全一致。

\section{常见误区}

\textbf{误区 1：}「两个平衡点都在定义域内」

错误。$x$ 的定义域是 $[x_0, +\infty)$（回收量从初值出发且单调递增），
$x = 0$ 远在定义域外，系统永远无法到达。物理上只存在一个平衡点 $x^* = K_{\text{eff}}$。

\medskip
\textbf{误区 2：}「不用逐因子分析，直接代个数就行」

可以，但必须\textbf{每个区间各代一个数}。不能只代一个数就下结论。
如果只代 $x = 2K_{\text{eff}}$（区间 II），得到 $f(2K_{\text{eff}}) < 0$，
就错误地认为``整个区间内 $f(x) < 0$''——这忽略了区间 I 中 $f(x) > 0$ 的事实。

\medskip
\textbf{误区 3：}「$f(x) < 0$ 就是系统在衰退」

不完全是。$f(x) < 0$ 只说明 $x$ 在减小，但如果 $x > K_{\text{eff}}$，
$x$ 减小恰恰是在\textbf{回到}平衡值——这是稳定性的体现，不是衰退。

\end{document}
